\subsubsection{Engajamento por Câmera}
O Engajamento por Câmera (\textit{Camera Engagement}) mede a proporção do tempo de sessão durante a qual os alunos mantiveram sua câmera de vídeo ativa. Esta métrica serve como indicador de participação ativa e presença social em ambientes de aprendizagem remota, onde a visibilidade facial pode influenciar a qualidade da interação e o senso de comunidade.

O cálculo é realizado pela função \textbf{calculate\_camera\_engagement(traces)} que implementa a seguinte fórmula:

\begin{equation}
\text{Engajamento por Câmera} = \frac{\text{Tempo com Câmera Ativa}}{\text{Tempo Total de Sessão}}
\end{equation}

\lstinputlisting[
    language=python,
    caption={Cálculo do Engajamento por Câmera},
    firstline=8,
    lastline=66,
    label={lst:calculate_camera_engagement}
]{scripts/metricsCalc/participation.py}


\textbf{Etapas do cálculo}:
\begin{enumerate}
    \item \textbf{Agrupamento por aluno}: Organiza os \textit{traces} por usuário
    \item \textbf{Ordenação temporal}: Ordena cronologicamente os registros de cada aluno
    \item \textbf{Cálculo incremental}: Para cada intervalo entre \textit{traces} consecutivos:
    \begin{itemize}
        \item Calcula a diferença de tempo em minutos
        \item Adiciona ao tempo total de sessão
        \item Se a câmera estava ativa no \textit{trace} anterior, adiciona ao tempo de câmera ativa
    \end{itemize}
    \item \textbf{Razão individual}: Calcula proporção de tempo com câmera ativa para cada aluno
    \item \textbf{Média agregada}: Computa média aritmética das proporções individuais
\end{enumerate}

\textbf{Detecção de câmera ativa}:
\begin{itemize}
    \item \texttt{cameraEnabled}: Indica se a câmera está habilitada no sistema
    \item \texttt{cameraStreaming}: Indica se a câmera está transmitindo vídeo ativamente
    \item \textbf{Lógica OR}: Considera câmera ativa se qualquer um dos indicadores for verdadeiro
\end{itemize}
